\documentclass[11pt,a4paper]{article}

\usepackage[T1]{fontenc}
\usepackage[utf8]{inputenc}
\usepackage[ngerman]{babel}
\usepackage[a4paper,margin=14mm]{geometry}
\usepackage{lmodern}
\usepackage{xcolor}
\usepackage{tabularx}
\usepackage{array}
\usepackage{microtype}
\usepackage{tikz}
\usetikzlibrary{arrows.meta,shapes.geometric}
\usepackage{textcomp}

% Gerade Anführungszeichen; " ist bei ngerman ein Kurzbefehl.
\newcommand{\qd}[1]{\textquotedbl#1\textquotedbl}

\pagestyle{empty}
\setlength{\parindent}{0pt}
\setlength{\parskip}{0pt}
\renewcommand{\familydefault}{\sfdefault}
\linespread{1.3}\selectfont

\definecolor{Navy}{HTML}{17324D}
\definecolor{Blue}{HTML}{3B6EA5}
\definecolor{Sky}{HTML}{EAF3FA}
\definecolor{Mint}{HTML}{E8F5EF}
\definecolor{Gold}{HTML}{F3B33D}
\definecolor{Ink}{HTML}{1E2935}
\definecolor{Muted}{HTML}{617080}
\definecolor{Line}{HTML}{D8E1E8}

\color{Ink}

\newcommand{\sectiontitle}[2]{%
  \vspace{4mm}
  {\color{#1}\bfseries\large #2}\par
  \vspace{1.2mm}
  {\color{#1}\rule{\linewidth}{0.7pt}}\par
  \vspace{2mm}
}

\newcommand{\footerline}{%
  {\color{Line}\rule{\linewidth}{0.5pt}}\par
  \vspace{1.5mm}
  \begin{tabularx}{\linewidth}{@{}Xr@{}}
    {\small\color{Muted}Klasse 9 · Tabellenkalkulation} &
    {\small\color{Muted}Sicherungsblatt · Aufgabe 5}
  \end{tabularx}%
}

\begin{document}

\colorbox{Navy}{%
  \parbox{\dimexpr\linewidth-2\fboxsep\relax}{%
    \vspace{3mm}
    \color{white}
    \begin{tabularx}{\linewidth}{@{}Xr@{}}
      {\small\bfseries INFORMATIK · KLASSE 9} &
      \colorbox{Gold}{\color{Navy}\small\bfseries\strut LÖSUNGEN}
    \end{tabularx}
    \vspace{2.5mm}

    {\fontsize{22}{25}\selectfont\bfseries Sicherungsblatt: Aufgabe 5}\par
    \vspace{1mm}
    {\large Datenflussdiagramme}
    \vspace{3mm}
  }%
}

\vspace{4mm}
\colorbox{Sky}{%
  \parbox{\dimexpr\linewidth-2\fboxsep\relax}{%
    \vspace{2.4mm}
    \textcolor{Blue}{\bfseries Ziel:}
    Du kannst ein \textbf{Datenflussdiagramm} lesen, selbst eines zeichnen,
    die WENN-Funktion als Diagramm darstellen und einen Fehler darin finden.
    \vspace{2.4mm}
  }%
}

\sectiontitle{Blue}{1 · Bausteine}

\begin{tikzpicture}[x=1mm,y=1mm,
    term/.style={font=\bfseries\small, text=Blue, anchor=north west,
                 inner sep=0pt},
    desc/.style={font=\small\linespread{1.0}\selectfont, text=Ink,
                 anchor=north west, inner sep=0pt,
                 text width=64mm, align=left,
                 execute at begin node={\hyphenpenalty=10000\relax}},
    func/.style={draw=Ink, line width=0.7pt, rounded corners=2pt, fill=white,
                 minimum width=20mm, minimum height=7mm, inner sep=1pt,
                 font=\small},
    io/.style={font=\small, inner sep=1pt},
    formula/.style={font=\small, text=Blue, inner sep=1pt},
    flow/.style={draw=Ink, line width=0.7pt,
                 -{Stealth[length=2.2mm,width=1.8mm]}},
    wire/.style={draw=Ink, line width=0.7pt},
    pointer/.style={draw=Blue, line width=0.5pt,
                    dash pattern=on 1.2mm off 1.0mm,
                    shorten <=1.5mm, shorten >=1.5mm,
                    -{Stealth[length=1.8mm,width=1.5mm]}}]

  % ---- Begriffsspalte ----
  \node[term] (t-eingabe) at (0,102) {Eingabe};

  \node[term] (t-verteiler) at (0,89) {Verteiler};
  \node[desc] at (0,84.5) {Ein kleiner Kreis.\\
    Er gibt denselben Wert an mehrere Funktionen weiter.};

  \node[term] (t-konstante) at (0,65.5) {Konstante};

  \node[term] (t-funktion) at (0,56.5) {Funktion};
  \node[desc] at (0,52) {Ein Rechen- oder Vergleichsschritt.\\
    Sie hat mehrere Eingangspfeile, aber genau einen Ausgangspfeil.};

  \node[term] (t-ausgabe) at (0,33) {Ausgabe};

  % ---- Beispielhaftes Datenflussdiagramm ----
  \node[io] (preis) at (156,100) {Preis};
  \draw[wire] (156,97) -- (156,84);
  \draw[flow] (156,84) -- (156,33.5);
  \node[circle, fill=Ink, inner sep=0pt, minimum size=1.9mm]
    (vert) at (156,84) {};
  \draw[wire] (156,84) -- (138,84);
  \draw[flow] (138,84) -- (138,66.5);

  \node[io] (konst) at (126,74) {0,1};
  \draw[flow] (126,71) -- (126,66.5);

  \node[func] (mal) at (132,63) {mal};
  \draw[wire] (132,59.5) -- (132,48) -- (146,48);
  \draw[flow] (146,48) -- (146,33.5);

  \node[func] (plus) at (151,30) {plus};
  \draw[flow] (151,26.5) -- (151,15.5);
  \node[io] (endpreis) at (151,12) {Endpreis};

  \node[formula] at (140,2) {Preis $\cdot$ 0,1 + Preis = Endpreis};

  % ---- Zeigerpfeile ----
  \draw[pointer] (t-eingabe.east) -- (preis.west);
  \draw[pointer] (t-verteiler.east)
    .. controls (70,92) and (132,93) .. (vert.center);
  \draw[pointer] (t-konstante.east) -- (konst.west);
  \draw[pointer] (t-funktion.east) -- (mal.west);
  \draw[pointer] (t-ausgabe.east) -- (endpreis.west);

\end{tikzpicture}

\vspace{7mm}
\colorbox{Mint}{%
  \parbox{\dimexpr\linewidth-2\fboxsep\relax}{%
    \vspace{2.4mm}
    \textcolor{Blue}{\bfseries Merke:}\\[2.4mm]
    Ein \textbf{Datenflussdiagramm} zeigt, wie Werte von den \textbf{Eingaben}
    durch \textbf{Funktionen} bis zu den \textbf{Ausgaben} fließen.\\[2.4mm]
    Die Pfeile geben die \textbf{Richtung des Datenflusses} an.\\[2.4mm]
    Ein Verteiler gibt denselben Wert an mehrere Funktionen weiter.\\[2.4mm]
    Eine Konstante ist ein fest eingetragener Wert, der nicht eingegeben
    wird.\\[2.4mm]
    Jede Funktion entspricht einem Teilschritt einer Formel.\\[2.4mm]
    Die \textbf{WENN-Funktion} hat drei Eingänge: Wahrheitswert, Ja-Fall,
    Nein-Fall.
    \vspace{2.4mm}
  }%
}

\vfill
\footerline

\newpage

\sectiontitle{Blue}{2 · Aussagefunktion}

\begin{tikzpicture}[x=1mm,y=1mm,
    desc/.style={font=\small\linespread{1.15}\selectfont, text=Ink,
                 anchor=north west, inner sep=0pt,
                 text width=104mm, align=left},
    func/.style={draw=Ink, line width=0.7pt, ellipse, fill=white,
                 minimum width=22mm, minimum height=8mm, inner sep=1pt,
                 font=\small},
    io/.style={font=\small, inner sep=1pt},
    hint/.style={font=\scriptsize\bfseries, text=Blue, inner sep=1pt},
    flow/.style={draw=Ink, line width=0.7pt,
                 -{Stealth[length=2.2mm,width=1.8mm]}}]

  \node[io] (note) at (18,32) {Note};
  \node[io] (zwei) at (46,32) {2};
  \node[func] (kleiner) at (32,19) {kleiner};
  \draw[flow] (note.south) -- (kleiner.north west);
  \draw[flow] (zwei.south) -- (kleiner.north east);
  \node[io] (wert) at (32,4) {Wahrheitswert};
  \draw[flow] (kleiner.south) -- (wert.north);
  \node[hint] at (32,0) {WAHR oder FALSCH};

  \node[desc] at (76,33) {Eine \textbf{Aussagefunktion} vergleicht Werte und
    beantwortet eine Ja-Nein-Frage. Ihr Ergebnis ist ein
    \textbf{Wahrheitswert}: WAHR oder FALSCH.\\[1.5mm]
    Note 1: $1 < 2$ ist WAHR.\quad Note 2: $2 < 2$ ist FALSCH.\\[1.5mm]
    Weitere Aussagefunktionen: größer, größer gleich, gleich. ISTGERADE hat
    nur einen Eingang: Aus 12 wird WAHR.};
\end{tikzpicture}

\sectiontitle{Blue}{3 · WENN-Funktion als Datenflussdiagramm}

\begin{tikzpicture}[x=1mm,y=1mm,
    desc/.style={font=\small\linespread{1.15}\selectfont, text=Ink,
                 anchor=north west, inner sep=0pt,
                 text width=66mm, align=left},
    func/.style={draw=Ink, line width=0.7pt, ellipse, fill=white,
                 minimum width=22mm, minimum height=8mm, inner sep=1pt,
                 font=\small},
    io/.style={font=\small, inner sep=1pt},
    hint/.style={font=\scriptsize\bfseries, text=Blue, inner sep=1pt},
    flow/.style={draw=Ink, line width=0.7pt,
                 -{Stealth[length=2.2mm,width=1.8mm]}}]

  \node[io] (note) at (16,62) {Note};
  \node[io] (zwei) at (44,62) {2};
  \node[func] (kleiner) at (30,50) {kleiner};
  \draw[flow] (note.south) -- (kleiner.north west);
  \draw[flow] (zwei.south) -- (kleiner.north east);

  \node[hint] at (72,44) {2 · Ja-Fall};
  \node[io] (ja) at (72,39.5) {\qd{5\texteuro}};
  \node[hint] at (100,44) {3 · Nein-Fall};
  \node[io] (nein) at (100,39.5) {\qd{0\texteuro}};

  \node[func] (wenn) at (62,22) {WENN};
  \draw[flow] (kleiner.south) -- (wenn.north west);
  \node[hint, anchor=east] at (41,35) {1 · Wahrheitswert};
  \draw[flow] (ja.south) -- (wenn.north);
  \draw[flow] (nein.south) -- (wenn.north east);

  \node[io] (geld) at (62,6) {Notengeld};
  \draw[flow] (wenn.south) -- (geld.north);

  \node[desc] at (116,62) {Die \textbf{WENN-Funktion} hat drei Eingänge:\\[1mm]
    \textbf{1.}~den Wahrheitswert einer Aussagefunktion,\\
    \textbf{2.}~den Wert für den Ja-Fall,\\
    \textbf{3.}~den Wert für den Nein-Fall.\\[1.5mm]
    Bei WAHR gibt sie den Ja-Fall aus, bei FALSCH den Nein-Fall.\\[1.5mm]
    Note 1 \textrightarrow{} WAHR \textrightarrow{} \qd{5\texteuro}\\
    Note 3 \textrightarrow{} FALSCH \textrightarrow{} \qd{0\texteuro}};
\end{tikzpicture}

\vspace{3mm}
{\small\textcolor{Blue}{\bfseries Termdarstellung:}\quad
\texttt{=WENN(Note\textless 2; \qd{5\texteuro}; \qd{0\texteuro})}}\par

\vspace{7mm}
\colorbox{Mint}{%
  \parbox{\dimexpr\linewidth-2\fboxsep\relax}{%
    \vspace{2.4mm}
    \textcolor{Blue}{\bfseries Merke:}\\[2.4mm]
    Eine \textbf{Aussagefunktion} gibt einen \textbf{Wahrheitswert} aus: WAHR
    oder FALSCH.\\[2.4mm]
    Die \textbf{WENN-Funktion} hat drei Eingänge: den Wahrheitswert einer
    Aussagefunktion, den Wert für den Ja-Fall und den Wert für den
    Nein-Fall.\\[2.4mm]
    In der Termdarstellung stehen die drei Eingänge in genau dieser
    Reihenfolge, getrennt durch Semikolons.
    \vspace{2.4mm}
  }%
}

\vfill
\footerline

\end{document}
